% sec5_3.tex — Section 5.3: Unbalanced Panels (optional)

\section{Unbalanced Panels (optional)}
\label{sec:5.3}

In applying the multiple-equation model to panel data, we have been making an
implicit but important assumption that the variables are observable for all $m$, so
that the number of observations for each cross-section unit is the same ($M$). Such
a panel is called a \textbf{balanced panel}. But no available panel data sets are balanced,
thanks to exits from and entries to the survey. Inevitably, some firms disappear
from the sample due to bankruptcies and mergers before the end year $M$ of the
survey, while those firms that did not exist at the start of the survey may be included
later. Some of the households initially in the survey will not complete the spell of
$M$ periods due to household resolutions or because respondents get tired of being
asked the same questions over and over again. The panel of identical siblings is
unbalanced if it includes identical triplets (and even septuplets) as well as twins.

This section shows that the fixed-effects estimator easily extends to unbalanced
panels. This sanguine conclusion, however, depends on the assumption that
whether the cross-section unit stays in the sample does not depend on the error
term. Without this assumption, the estimator is no longer consistent. This phenomenon
is called the \textbf{selectivity bias}. A full discussion of why a selectivity bias
arises and how it can be corrected for is relegated to Section~8.2, because the issue
is not specific to panel data.

\subsection*{``Zeroing Out'' Missing Observations}

To handle missing observations, it is convenient to define a dummy variable, $d_{im}$,
\begin{equation}
  d_{im} = \begin{cases}
    1 & \text{if observation $m$ is in the sample,} \\
    0 & \text{otherwise,}
  \end{cases}
  \label{eq:5.3.1} \tag{5.3.1}
\end{equation}
and
\begin{equation}
  \underset{(M \times 1)}{\mathbf{d}_i}
  = \begin{bmatrix} d_{i1} \\ \vdots \\ d_{iM} \end{bmatrix}, \quad
  M_i = \sum_{m=1}^{M} d_{im} = \text{\# observations from } i.
  \label{eq:5.3.2} \tag{5.3.2}
\end{equation}

If observation $m$ is missing for group $i$, fill the $m$-th rows of $\mathbf{y}_i$, $\mathbf{F}_i$, and
$\boldsymbol{\eta}_i$ with zeros so that they continue to be of fixed size. That is,
\begin{equation}
  \underset{(M \times 1)}{\mathbf{y}_i}
  = \begin{bmatrix} d_{i1} \cdot y_{i1} \\ \vdots \\ d_{iM} \cdot y_{iM} \end{bmatrix}, \quad
  \underset{(M \times \#\boldsymbol{\beta})}{\mathbf{F}_i}
  = \begin{bmatrix} d_{i1} \cdot \mathbf{f}_{i1}' \\ \vdots \\ d_{iM} \cdot \mathbf{f}_{iM}' \end{bmatrix}, \quad
  \underset{(M \times 1)}{\boldsymbol{\eta}_i}
  = \begin{bmatrix} d_{i1} \cdot \eta_{i1} \\ \vdots \\ d_{iM} \cdot \eta_{iM} \end{bmatrix}.
  \label{eq:5.3.3} \tag{5.3.3}
\end{equation}

Thus, for each $i$ and $m$, either all the elements of the vector
$(y_{im}, \mathbf{f}_{im})$ are observable
or none is observable. We will not consider the case where some, but not all,
elements of $(y_{im}, \mathbf{f}_{im})$ are observable.

With this new notation, \eqref{eq:5.1.1pp} on page~329 becomes
\begin{equation}
  \mathbf{y}_i = \mathbf{F}_i\boldsymbol{\beta} + \mathbf{d}_i \cdot \mathbf{b}_i'\boldsymbol{\gamma}
  + \mathbf{d}_i \cdot \alpha_i + \boldsymbol{\eta}_i
  \quad (i = 1, 2, \ldots, n).
  \label{eq:5.3.4} \tag{5.3.4}
\end{equation}

In the case of balanced panels, where $\mathbf{d}_i = \mathbf{1}_M$, the transformation matrix $\mathbf{Q}$ in the
fixed-effects formula \eqref{eq:5.2.4} is the annihilator associated with $\mathbf{1}_M$
(see \eqref{eq:5.1.10}).

With missing observations, we use as the transformation matrix the annihilator
associated with $\mathbf{d}_i$:
\begin{align}
  \underset{(M \times M)}{\mathbf{Q}}
  &= \mathbf{I}_M - \mathbf{d}_i(\mathbf{d}_i'\mathbf{d}_i)^{-1}\mathbf{d}_i'
  \notag \\
  &= \mathbf{I}_M - \frac{1}{M_i}\,\mathbf{d}_i\mathbf{d}_i'
  \quad (\text{since } \mathbf{d}_i'\mathbf{d}_i = M_i).
  \label{eq:5.3.5} \tag{5.3.5}
\end{align}

This matrix depends on $i$ because $\mathbf{d}_i$ differs across $i$. Nevertheless, we leave this
dependence of $\mathbf{Q}$ on $i$ implicit for economy of notation. With $\mathbf{y}_i$, $\mathbf{F}_i$,
$\boldsymbol{\eta}_i$, and $\mathbf{Q}$ thus
redefined, the formula for the fixed-effects estimator remains the same as \eqref{eq:5.2.4}.
Substituting \eqref{eq:5.3.4} into this formula and observing that $\mathbf{Q}\mathbf{d}_i = \mathbf{0}$, we see that the
expression for the sampling error, too, is the same as \eqref{eq:5.2.5}.

\subsection*{Zeroing Out versus Compression}

As seen in Section~5.2 for the case of balanced panels, the fixed-effects estimator
can be calculated as the OLS estimator on the pooled sample (of size $Mn$) of deviations
from group means. As before, let $\tilde{\mathbf{y}}_i$, $\widetilde{\mathbf{F}}_i$, and
$\tilde{\boldsymbol{\eta}}_i$ be the transformations by $\mathbf{Q}$
of $\mathbf{y}_i$, $\mathbf{F}_i$, and $\boldsymbol{\eta}_i$, respectively. For example,
$\tilde{\mathbf{y}}_i$ looks like
\[
  \underset{(M \times M)}{\tilde{\mathbf{y}}_i} \equiv \mathbf{Q}\mathbf{y}_i
  = \begin{bmatrix} d_{i1} \cdot y_{i1} - d_{i1} \cdot \bar{y}_i \\
    \vdots \\
    d_{iM} \cdot y_{iM} - d_{iM} \cdot \bar{y}_i
  \end{bmatrix},
\]
where
\begin{equation}
  \bar{y}_i \equiv \frac{1}{M_i}\,\mathbf{d}_i'\mathbf{y}_i
  = \frac{1}{M_i} \times (\text{sum of } y_{im} \text{ over } m
  \text{ for which data are available}).
  \label{eq:5.3.6} \tag{5.3.6}
\end{equation}

Thus, if observation $m$ for group $i$ is available, the $m$-th element of $\tilde{\mathbf{y}}_i$ is still the
deviation of $y_{im}$ from the group mean, but the group mean is over available observations.
Since $d_{im} = 0$ if observation $m$ is not available, the rows corresponding to
missing observations in $\tilde{\mathbf{y}}_i$ are zero. The same structure applies to columns of
$\widetilde{\mathbf{F}}_i$.
Therefore, the pooled sample of size $Mn$ created from
$(\tilde{\mathbf{y}}_i, \widetilde{\mathbf{F}}_i)$ $(i = 1, \ldots, n)$ has
a number of rows filled with zeros. Those rows can be dropped from the pooled
sample without affecting the numerical value of the fixed-effects estimator. That
is, we can ``compress'' $(\tilde{\mathbf{y}}_i, \widetilde{\mathbf{F}}_i)$ by dropping rows filled with zeros, before forming
the pooled sample.

\subsection*{No Selectivity Bias}

The fixed-effects estimator, thus adapted to unbalanced panels, remains consistent
and asymptotically normal, if the orthogonality conditions \eqref{eq:5.2.6} in Proposition
5.1 are modified to
\begin{equation}
  \text{(no selectivity bias)} \quad
  \E(\mathbf{f}_{im} \cdot \eta_{ih} \mid \mathbf{d}_i) = \mathbf{0}
  \quad (m, h = 1, 2, \ldots, M).
  \label{eq:5.3.7} \tag{5.3.7}
\end{equation}

As shown in the next proposition, this condition is crucial for ensuring the consistency
of the fixed-effects estimator. By the Law of Total Expectations, this condition
is stronger than \eqref{eq:5.2.6} that the unconditional expectation is zero. It is easy to
construct an example where \eqref{eq:5.3.7} does not hold because the pattern of selection
$\mathbf{d}_i$ depends on the error term $\boldsymbol{\eta}_i$. Note that \eqref{eq:5.3.7} does not involve the fixed effect
$\alpha_i$. Thus, if the dependence of sample selection on the error term is only through
the fixed effect, the problem of sample selectivity bias does not occur.

If the orthogonality conditions \eqref{eq:5.2.6} are strengthened as \eqref{eq:5.3.7}, the conclusions
of Proposition~5.3 carry over:

\begin{proposition}[large-sample properties of the fixed-effects estimator for unbalanced panels]
\label{prop:5.4}
Consider the error-components model consisting of \eqref{eq:5.3.4},
\eqref{eq:5.1.2}, \eqref{eq:5.1.4}--\eqref{eq:5.1.6}, \eqref{eq:5.1.15}, and
\eqref{eq:5.3.7}, with $(\mathbf{y}_i, \mathbf{F}_i, \boldsymbol{\eta}_i)$ as in \eqref{eq:5.3.3}. As
before, define $(\tilde{\mathbf{y}}_i, \widetilde{\mathbf{F}}_i, \tilde{\boldsymbol{\eta}}_i)$
by \eqref{eq:5.2.2}, but with $\mathbf{Q}$ given by \eqref{eq:5.3.5}. Then the conclusions
of Proposition~5.3 hold.
\end{proposition}

Here we show that $\E(\mathbf{F}_i'\mathbf{Q}\boldsymbol{\eta}_i) = \mathbf{0}$.
(The other parts of the proof, which are applications of Kolmogorov's LLN and the Lindeberg-Levy CLT, are the same as in the
proof of Proposition~5.3.) Since
\begin{equation}
  \mathbf{F}_i'\mathbf{Q}\boldsymbol{\eta}_i
  = \sum_{m=1}^{M}\sum_{h=1}^{M}q_{mh}^{(i)} \cdot d_{im} \cdot d_{ih}
  \cdot \mathbf{f}_{im} \cdot \eta_{ih},
  \label{eq:5.3.8} \tag{5.3.8}
\end{equation}
where $q_{mh}^{(i)}$ is the $(m,h)$ element of $\mathbf{Q}$ applied to the $i$-th equations, we have
\begin{equation}
  \E(\mathbf{F}_i'\mathbf{Q}\boldsymbol{\eta}_i)
  = \sum_{m=1}^{M}\sum_{h=1}^{M}
  \E(q_{mh}^{(i)} \cdot d_{im} \cdot d_{ih} \cdot \mathbf{f}_{im} \cdot \eta_{ih}).
  \label{eq:5.3.9} \tag{5.3.9}
\end{equation}

If there is no selectivity bias so that
$\E(\mathbf{f}_{im} \cdot \eta_{ih} \mid \mathbf{d}_i) = \mathbf{0}$, then \eqref{eq:5.3.9} equals zero
because
\begin{align}
  &\E(q_{mh}^{(i)} \cdot d_{im} \cdot d_{ih} \cdot \mathbf{f}_{im} \cdot \eta_{ih})
  \notag \\
  &\quad= \E\!\bigl[q_{mh}^{(i)} \cdot d_{im} \cdot d_{ih}
  \cdot \E(\mathbf{f}_{im} \cdot \eta_{ih} \mid \mathbf{d}_i)\bigr]
  \quad (\text{since } q_{mh}^{(i)} \text{ is a function of } \mathbf{d}_i).
  \label{eq:5.3.10} \tag{5.3.10}
\end{align}

It is interesting that the extension to unbalanced panels is easier without conditional
homoskedasticity. In order for \eqref{eq:5.2.7} of Proposition~5.1 to be the asymptotic
variance of the estimator under conditional homoskedasticity, we would have to
assume that the second moment of $\tilde{\boldsymbol{\eta}}_i$ conditional on $\mathbf{x}_i$ \emph{and}
$\mathbf{d}_i$ equals the unconditional second moment, and consistent estimation of this second moment would be
a bit complicated.

\bigskip
\noindent\rule{\textwidth}{0.5pt}
\textsc{Questions for Review}
\medskip

\begin{enumerate}
\item For $\mathbf{d}_i = (1, 0, 1)'$, write down $\mathbf{Q}$.

\item (Alternative derivation of the fixed-effects estimator for unbalanced panels)\quad
  Let $\mathbf{y}_i^*$ be the $M_i$-dimensional vector created from the $M$-dimensional vector
  $\mathbf{y}_i$ by dropping rows for which there is no observation. Similarly, create $\mathbf{F}_i^*$
  ($M_i \times \#\boldsymbol{\beta}$) from $\mathbf{F}_i$ and $\mathbf{d}_i^*$ from $\mathbf{d}_i$
  (so $\mathbf{d}_i^*$ is the $M_i$-dimensional vector of
  ones). Let $\mathbf{Q}^*$ ($M_i \times M_i$) be the annihilator associated with $\mathbf{d}_i^*$.
  Verify that the
  fixed-effects estimator can also be written as \eqref{eq:5.2.4} in terms of those starred
  matrices and vectors.

\item Section~4.6 described the pooled OLS in levels for the system (4.6.1$'$) on page~292.
  \begin{enumerate}[label=(\alph*)]
  \item How would you modify the estimator (4.6.11$'$) on page~293 to accommodate
    unbalanced panels? \textbf{Hint:} Define $\mathbf{y}_i$ and $\mathbf{Z}_i$
    ($M \times L$) by ``zeroing out.''
  \item Without conditional homoskedasticity, what is its asymptotic variance?
    [Answer: $\bigl[\E(\mathbf{Z}_i\mathbf{Z}_i')\bigr]^{-1}\,
    \E(\mathbf{Z}_i\boldsymbol{\varepsilon}_i\boldsymbol{\varepsilon}_i'\mathbf{Z}_i')\,
    \bigl[\E(\mathbf{Z}_i\mathbf{Z}_i')\bigr]^{-1}$.]
  \item How would you estimate the asymptotic variance? \textbf{Hint:} Replace
    $\widetilde{\mathbf{F}}_i$ by $\mathbf{Z}_i$,
    $\mathbf{Q}$ by $\mathbf{I}_M$, and $\tilde{\boldsymbol{\eta}}_i$ by
    $\boldsymbol{\varepsilon}_i$ in Proposition~5.3.
  \end{enumerate}
\end{enumerate}
